% Generated draft: requires mathematical and source review.
\documentclass[11pt]{article}
\usepackage[margin=25mm]{geometry}
\usepackage{fontspec}
\setmainfont{Latin Modern Roman}
\usepackage{amsmath,amssymb,mathtools,unicode-math}
\setmathfont{Latin Modern Math}
\usepackage{enumitem,needspace,xurl,hyperref}
\setlength{\parindent}{0pt}
\setlength{\parskip}{5pt}
\setlength{\emergencystretch}{3em}
\newcommand{\N}{\mathbb{N}}
\newcommand{\Z}{\mathbb{Z}}
\newcommand{\Q}{\mathbb{Q}}
\newcommand{\R}{\mathbb{R}}
\newcommand{\C}{\mathbb{C}}
\newcommand{\F}{\mathbb{F}}
\newcommand{\A}{\mathbb{A}}
\newcommand{\PP}{\mathbb{P}}
\newcommand{\E}{\mathbb{E}}
\newcommand{\eps}{\varepsilon}
\newcommand{\dd}{\,\mathrm{d}}
\begin{document}
\section*{UnsolvedMath ID: 800358}
\textbf{Source problem code:} catalogue entry
\subsection{Definitions and mathematical statement}
Let \(D=\{n\in\Z:n\geq4\}\), equipped with the topology induced by the \(2\)-adic metric. For integers \(a\neq b\), this metric is \(d_2(a,b)=2^{-v_2(a-b)}\), where \(v_2(t)\) is the largest nonnegative integer \(e\) such that \(2^e\mid t\); set \(d_2(a,a)=0\). An integer-valued function \(F:D\to\Z\), with the same metric on its codomain, is \(2\)-adically continuous precisely when
\[
\forall n_0\in D\;\forall k\in\Z_{\geq1}\;\exists \ell\in\Z_{\geq0}\;\forall n\in D,\qquad
n\equiv n_0\pmod{2^\ell}\ \Longrightarrow\ F(n)\equiv F(n_0)\pmod{2^k}.
\]
Here \(\ell\) may depend on both \(n_0\) and \(k\).

The supplied question is whether the sequence denoted in the source by
\[
F(n)=T_{113}(n,2n-3;3),\qquad n\in D,
\]
is \(2\)-adically continuous in this sense. The supplied record does not define \(T_{113}(a,b;c)\), the subscript \(113\), the underlying benzel regions, the admissible trimer tilings, or the original weighting. Consequently, the displayed continuity condition gives a precise target for an integer-valued sequence once the intended enumeration is supplied, but the record alone does not specify that sequence completely. No replacement definition of \(T_{113}\) is asserted here.

The comment immediately following the question reports apparent constancy modulo \(2\) and modulo \(4\), apparent period \(2\) modulo \(8\), and apparent period \(8\) modulo \(16\), with repeating residue block
\[
(14,6,10,2,6,14,2,10).
\]
It states that the case \(n=3\) breaks the pattern and refers to OEIS A364438. The starting index of this residue block is not specified. These are supplied experimental observations, not established congruences; the stated exception \(n=3\) lies outside the question's domain.

A separate comment preceding the question reports apparent constancy modulo \(2\), \(4\), and \(8\), and apparent period \(8\) modulo \(16\), with repeating residue block
\[
(4,4,12,12,4,12,12,4).
\]
It states that \(n=4\) breaks the pattern and refers to OEIS A364418. The sequence to which this preceding comment belongs, and the indexing of its residue block, are not identified in the supplied excerpt. It therefore cannot safely be attached to \(T_{113}(n,2n-3;3)\).

The source also states that, for a fixed benzel, the number of right stones minus the number of left stones is independent of the tiling. On this basis it describes the weighting in Problems 17 and 18, up to an overall power of \(3\), by assigning weight \(1\) to left stones and weight \(9\) to right stones. It further proposes assigning weights \(1\) and \(m\), respectively, and reports apparent \(p\)-adic continuity when \(m+1\) is a power of a prime \(p\), together with apparent \(5\)-adic continuity for \(m=9\). These broader observations are preserved as source qualifications: the permitted values of \(m\), the precise resulting enumerations, their parameter domains, and their intended continuity statements are not supplied. The author describes these patterns as possible subjects of future work.
\subsection{Short English statement}
Propp asks whether \(T_{113}(n,2n-3;3)\), for integers \(n\geq4\), is continuous in the \(2\)-adic topology: sufficiently strong congruence of inputs near each fixed input should imply any prescribed congruence of outputs. The record omits the definition of the enumeration, so the statement remains conditional on recovering that definition. Its congruence patterns and broader weighted observations are reported evidence, not proved claims.
\subsection{Sources}
\begin{itemize}
\item Local UnsolvedMath registry, permanent ID \(800358\), exact source problem code \texttt{catalogue entry}, title ``Trimer Covers in the Triangular Grid --- Problem 18.'' The record attributes the problem to Propp, proposed in 2022, and supplies the status label ``open''; that label is not an independently verified mathematical or literature assessment. Upstream revision: \texttt{d371ee917ac9647fbb06cd6fc529deec6f3dda82}.
\item Propp, \emph{Trimer Covers in the Triangular Grid: Twenty Mostly Open Problems} (2022), source item ``Problem 18 (\#10),'' \url{https://arxiv.org/abs/2206.06472}. The supplied provenance records access on 2026-10-16 and extraction by \texttt{arxiv\_latex}; these are supplied provenance statements, not a claim of independent source access or verification.
\item OEIS A364438, \url{https://oeis.org/A364438}, cited by the comment immediately following the continuity question.
\item OEIS A364418, \url{https://oeis.org/A364418}, cited by the preceding comment whose sequence is not identified in the supplied excerpt.
\item The supplied source attributes the alternative weighting and generalized continuity observations to the author during the final preparation of the article for publication, and records the author's intention to address such patterns in future work. No additional supporting reference is supplied for those observations.
\end{itemize}
\end{document}
